\chapter{Fundamentos clínicos y técnicos}
\label{chap:fundamentos}

\section{El electrocardiograma y sus componentes}
\label{sec:componentes_ecg}

El electrocardiograma registra las diferencias de potencial producidas por la actividad eléctrica del corazón mediante derivaciones que observan distintas proyecciones de esa actividad tridimensional. En el ECG convencional, diez electrodos permiten obtener doce derivaciones, de las cuales I, II, III, aVR, aVL y aVF describen el plano frontal, mientras que V1 a V6 corresponden al horizontal.

Para interpretar sus amplitudes y duraciones es necesario conocer la colocación de los electrodos, la ganancia, la velocidad de registro y el filtrado \cite{kligfield2007}. La representación habitual utiliza 25 mm por segundo y 10 mm por mV, de modo que un milímetro equivale a 40 ms en horizontal y 0,1 mV en vertical. Esta escala, visible en el formato clínico de la Figura~\ref{fig:signal_to_image}, puede perderse al recortar o redimensionar la imagen.

\begin{figure}[H]
  \centering
  \safeincludegraphics[width=0.78\textwidth]{sources/ptbxl_real_12lead_ecg.pdf}
  \caption{Representación clínica de un ECG real de doce derivaciones a partir del renderizado de un registro de PTB-XL/PhysioNet \cite{wagner2020}.}
  \label{fig:signal_to_image}
\end{figure}

La despolarización es la activación eléctrica de las células cardíacas y la repolarización, la recuperación de su estado eléctrico tras esa activación. La onda P representa la despolarización auricular y el intervalo PR, la propagación hacia los ventrículos. El complejo QRS corresponde a la despolarización ventricular. El punto J marca su final y el inicio del segmento ST. La onda T representa la repolarización ventricular. La Figura~\ref{fig:morfologia_ecg} sitúa estos componentes.

\begin{figure}[H]
  \centering
  \safeincludegraphics[width=0.86\textwidth]{results/fig_morfologia_ecg.pdf}
  \caption{Componentes morfológicos de un latido ECG. La figura identifica las ondas P, Q, R, S y T, el punto J y el segmento ST.}
  \label{fig:morfologia_ecg}
\end{figure}

Q es la primera deflexión negativa anterior a una positiva, R la primera positiva y S la negativa posterior a R. Una segunda deflexión positiva se denomina R prima. Estas letras describen la forma del trazado.

\section{Síndrome de Brugada y rasgos estudiados}
\label{sec:brugada_clinica}

El síndrome de Brugada es una canalopatía hereditaria, una alteración de los canales que regulan el paso de iones en las células cardíacas. Se relaciona con los canales de sodio y con variantes en el gen \emph{SCN5A}. Predispone a fibrilación ventricular y muerte súbita, especialmente durante el reposo o el sueño en varones jóvenes, aunque puede manifestarse a cualquier edad \cite{zeppenfeld2022,wilde2023scbt}. Las estimaciones citadas sitúan su prevalencia entre 1 y 5 casos por cada 10\,000 personas \cite{sieira2017}, con intervalos también descritos de una persona por cada 2\,000 a 5\,000. Puede estar infradiagnosticado porque muchos afectados permanecen asintomáticos.

El patrón tipo~1, o \emph{coved}, presenta elevación del ST de al menos 2 mm en V1 o V2, de forma convexa y descendente, seguida de T negativa. Los patrones tipo~2 y tipo~3 no bastan para establecer el diagnóstico \cite{zeppenfeld2022}. El patrón puede ser intermitente, por lo que un ECG basal normal no descarta el síndrome.

\subsection{Evaluación clínica}
\label{sec:evaluacion_clinica}

La interpretación combina el ECG con síntomas, antecedentes familiares y la exclusión de otras causas. El síncope inexplicado, la respiración agónica nocturna, la muerte súbita familiar y las arritmias ventriculares documentadas orientan la valoración del riesgo. El circuito comprende los siguientes pasos.

\begin{enumerate}
  \item \textbf{ECG basal de doce derivaciones.} Un patrón tipo~1 espontáneo permite establecer el diagnóstico tras descartar otras causas de elevación del ST.
  \item \textbf{Registro precordial alto.} Si el ECG es sospechoso, se registran V1 y V2 en el segundo o tercer espacio intercostal, en lugar del cuarto. La proximidad al tracto de salida del ventrículo derecho aumenta la sensibilidad \cite{brugadadrugsdiagnostic}.
  \item \textbf{Revisión clínica.} Se revisan síntomas, antecedentes, comorbilidades y fármacos, ya que los trastornos electrolíticos, la isquemia y algunos medicamentos pueden producir patrones similares. Durante un episodio de fiebre puede aparecer un patrón de Brugada que no se observa a temperatura normal \cite{zeppenfeld2022}.
  \item \textbf{Prueba farmacológica.} En pacientes seleccionados con sospecha persistente y ECG no diagnóstico se utilizan bloqueantes del canal de sodio, como ajmalina, flecainida, procainamida o pilsicainida. Se realiza con monitorización, personal entrenado y desfibrilador, y se considera positiva si aparece un patrón tipo~1 durante o después de administrar el fármaco \cite{wilde2023scbt}.
  \item \textbf{Estudio genético y decisión especializada.} En pacientes índice con patrón tipo~1 puede recomendarse estudiar \emph{SCN5A} para la caracterización molecular y el cribado familiar. La estratificación del riesgo y el tratamiento dependen de la valoración conjunta de especialistas.
\end{enumerate}

La Figura~\ref{fig:brugada_clinical_examples} muestra ejemplos clínicos de estos patrones.

\begin{figure}[H]
  \centering
  \begin{subfigure}[t]{0.58\textwidth}
    \centering
    \safeincludegraphics[width=\linewidth]{sources/brugada_type1_ecgpedia.pdf}
    \caption{ECG con morfología tipo 1 en V1 y V2.}
  \end{subfigure}
  \hfill
  \begin{subfigure}[t]{0.34\textwidth}
    \centering
    \safeincludegraphics[width=\linewidth]{sources/brugada_patterns_orphanet.pdf}
    \caption{Patrones tipo 1--3.}
  \end{subfigure}
  \caption{Ejemplos clínicos de morfología asociada al síndrome de Brugada. La imagen izquierda procede de ECGpedia, CC BY-SA 3.0 \cite{ecgpediaType1Image}. La imagen derecha procede de Napolitano y Priori, CC BY 2.0 \cite{napolitano2006}.}
  \label{fig:brugada_clinical_examples}
\end{figure}

\subsection{Morfología compatible con bloqueo de rama derecha}

El bloqueo de rama derecha (RBBB, por sus siglas en inglés) es una alteración de la conducción eléctrica que retrasa la activación del ventrículo derecho. Como esa activación ocurre después que la del izquierdo, cambia la parte final del complejo QRS y puede aparecer una segunda deflexión positiva en las derivaciones que observan el lado derecho del corazón.

En el bloqueo completo, el ECG presenta un QRS de al menos 120 ms, patrones como rsR prima o rSR prima en V1 o V2 y una onda S ensanchada en derivaciones laterales \cite{surawicz2009}. Estos criterios describen un problema de conducción eléctrica y algunos de sus rasgos visuales también aparecen en patrones asociados al síndrome de Brugada.

Como una forma semejante en una sola derivación no basta para diagnosticar el bloqueo, también deben valorarse la duración global del QRS, las derivaciones I y V6, la posición de los electrodos y otras causas de retraso terminal. Por ello, en esta memoria RBBB identifica únicamente un rasgo visual compatible, que también puede recordar al patrón tipo~1 de Brugada. A diferencia del bloqueo, este patrón concentra su onda J en las precordiales derechas y no requiere una onda S lateral recíproca \cite{corrado2010}. La Figura~\ref{fig:brugada_pattern} resume los rasgos considerados.

\begin{figure}[H]
  \centering
  \safeincludegraphics[width=0.94\textwidth]{results/fig_brugada_pattern.pdf}
  \caption{Representación esquemática de rasgos morfológicos asociados al patrón tipo 1 del síndrome de Brugada.}
  \label{fig:brugada_pattern}
\end{figure}

\subsection{Elevación del ST e inversión de T}

El ST se extiende desde J hasta el inicio de T. Su desplazamiento se mide respecto de la línea isoeléctrica, aproximada por TP o PR según el registro. La interpretación depende de la derivación, la magnitud y forma de la elevación, su distribución, la edad, el sexo y el cuadro clínico \cite{rautaharju2009st}. Puede aparecer en isquemia aguda, pericarditis, repolarización precoz, alteraciones de conducción y síndromes de onda J.

La polaridad de T varía entre derivaciones y con la edad y el sexo. También puede cambiar de forma secundaria a una despolarización anómala, con alteraciones discordantes del ST y de T. Una T negativa aislada no permite determinar si es primaria, secundaria, fisiológica o patológica.

La elevación del ST o la inversión de T pueden deberse a distintas causas, por lo que distinguir Brugada de otras enfermedades exige valorar el ECG completo junto con la información clínica. Las etiquetas del simulador se limitan a describir esas formas visibles, sin asignarles por sí solas un diagnóstico.

\section{Señal digital e imagen}

La señal digital conserva las amplitudes ordenadas en el tiempo, de modo que pueden filtrarse, medirse intervalos y compararse derivaciones con una escala conocida. Al representarla como imagen se mantiene la forma del trazado, pero se incorporan también la cuadrícula, el grosor de la línea, los márgenes y la distribución sobre el documento. Estos elementos pueden influir en la clasificación si se repiten más en una clase que en otra, aunque no correspondan a una alteración cardíaca.

Por ello, cuando solo se dispone de un documento, hay que elegir entre analizar su imagen o reconstruir primero la señal mediante digitalización. La segunda opción permite reutilizar modelos temporales, aunque puede introducir errores si los trazos se solapan o la calidad es insuficiente, mientras que el análisis directo debe distinguir las ondas de los elementos del fondo. La sección~\ref{sec:antecedentes_imagen} recoge los estudios que comparan estas alternativas.

\section{Clasificación y salida de los modelos}
\label{sec:clasificacion}

En clasificación multiclase las categorías son exclusivas. En clasificación multietiqueta pueden coexistir, como ocurre con los rasgos del ECG \cite{tsoumakas2007}. Para una imagen \(x\) y \(L\) etiquetas se define
\begin{equation}
\mathbf{y}=(y_1,\ldots,y_L), \qquad y_j\in\{0,1\}.
\end{equation}

La red produce un logit \(z_j\) por etiqueta, una puntuación numérica que puede tomar cualquier valor real. La función sigmoide lo transforma en una probabilidad entre cero y uno mediante la expresión
\begin{equation}
\hat{p}_j=\frac{1}{1+\exp(-z_j)}.
\end{equation}

La entropía cruzada binaria (BCE) es la función de pérdida que mide la discrepancia entre las probabilidades predichas y las etiquetas. Se utiliza para ajustar el modelo y se calcula como
\begin{equation}
\mathcal{L}_{\mathrm{BCE}}=
-\frac{1}{L}\sum_{j=1}^{L}
\left[y_j\log(\hat{p}_j)+(1-y_j)\log(1-\hat{p}_j)\right].
\label{eq:bce}
\end{equation}

Aunque la red produzca una salida por rasgo, todas se obtienen a partir de la misma imagen y pueden aprovechar asociaciones entre etiquetas. Si, por ejemplo, las T negativas apareciesen siempre junto con elevación del ST, el modelo podría predecir la primera a partir de la segunda sin examinar su región, por lo que acertar ambas etiquetas no demuestra que reconozca cada onda de forma independiente.

Para evaluar estas salidas, sus probabilidades se convierten en decisiones mediante un umbral elegido con los pacientes de validación y mantenido durante la prueba. Ajustarlo con los casos que se evalúan permitiría favorecer el resultado después de conocer sus respuestas, de ahí la separación descrita en la sección~\ref{sec:umbral_cnn}. En Gemma~4, la respuesta llega como texto estructurado y se transforma en las mismas etiquetas positivas o negativas para poder calcular las métricas de ambos métodos.

\section{Redes neuronales convolucionales}
\label{sec:fundamentos_cnn}

Una CNN combina capas convolucionales, que aprenden filtros para detectar formas, con funciones de activación, operaciones que reducen la resolución y una salida de clasificación. En este trabajo, esas capas procesan la imagen del ECG, donde las primeras pueden detectar bordes y cambios de dirección del trazo, mientras que las siguientes reúnen esa información para representar picos, valles y relaciones entre ondas. La convolución aplica un mismo filtro en distintas posiciones, por lo que no necesita aprender por separado una forma cada vez que aparece desplazada.

En una señal digital se utilizan filtros unidimensionales que, al apilarse las capas, abarcan intervalos temporales cada vez mayores. En la imagen, el recorrido es espacial, pero sigue siendo importante conservar la posición de las ondas, porque una deflexión negativa no significa lo mismo dentro del QRS que en T. Por ello, la tolerancia a pequeños desplazamientos debe combinarse con un preprocesamiento que alinee los latidos sin borrar sus relaciones temporales.

ResNet18 incorpora conexiones residuales que suman la entrada de un bloque a su transformación y permiten entrenar la red aprendiendo cambios sobre esa entrada \cite{he2016}. El modelo parte de pesos aprendidos con ImageNet y se comprueba experimentalmente si esos filtros resultan útiles para reconocer los trazos y texturas del ECG. La adaptación de la salida a las tareas del simulador y del HUCA se describe en la sección~\ref{sec:cnn_entrenamiento}.

\section{Explicabilidad y adaptación de dominio}
\label{sec:explicabilidad}

La inteligencia artificial explicable, o \textit{explainable AI} (XAI), estudia cómo examinar las decisiones de un modelo para que una persona pueda valorar en qué información se apoyan. En el análisis de ECG interesa comprobar si la clasificación se relaciona con las ondas relevantes o con elementos del documento, como la cuadrícula y los márgenes, antes de atribuirle una interpretación clínica.

Para inspeccionar la CNN se utiliza Grad-CAM, que construye un mapa de regiones asociadas a una predicción ponderando las activaciones de una capa convolucional con los gradientes de la clase \cite{selvaraju2017}. Así puede observarse si la respuesta se concentra alrededor del QRS, del ST o de T, aunque la resolución del mapa depende de la capa elegida y puede abarcar una zona amplia alrededor de un artefacto. Esta inspección ayuda a detectar comportamientos sospechosos, pero no permite asegurar por sí sola qué rasgo clínico determinó el resultado.

En un modelo de visión y lenguaje también puede solicitarse una explicación escrita, pero que resulte comprensible no garantiza que describa el proceso que produjo la predicción. Por eso se distinguen las etiquetas estructuradas utilizadas para evaluar el rendimiento de las justificaciones examinadas de forma cualitativa.

Esta precaución también importa al cambiar de conjunto de datos, porque la red puede aprender el estilo de representación en lugar de la forma de las ondas. La adaptación de dominio intenta reducir esa diferencia cuando se pasa de los latidos alineados y las cuadrículas uniformes del simulador a ECG con ruido, deriva de línea base, variación de frecuencia, anotaciones y trazos superpuestos. El capítulo~\ref{chap:simulador} describe los controles del generador y la sección~\ref{sec:adaptacion_cnn}, los métodos empleados para estudiar ese cambio.

\section{Modelos de visión y lenguaje y aprendizaje en contexto}
\label{sec:fundamentos_icl}

Un modelo de visión y lenguaje (VLM) combina un componente que procesa imágenes con un modelo generativo de texto. El componente visual transforma la imagen en vectores numéricos que representan su contenido y que el modelo relaciona con las instrucciones escritas. El entrenamiento multimodal permite aprender asociaciones entre imágenes y texto, de modo que una consulta puede pedir una descripción, una respuesta a una pregunta o una clasificación \cite{alayrac2022}.

Al analizar un ECG, la imagen aporta el trazado y las instrucciones indican qué debe examinarse y cómo expresar la respuesta. Así, el mismo modelo puede recibir preguntas sobre la forma de las ondas o responder con una etiqueta, según la tarea planteada.

\begin{samepage}
El aprendizaje en contexto (ICL) incorpora además ejemplos resueltos a la consulta. En su versión multimodal, cada demostración reúne una imagen y su etiqueta, con las que se muestra al modelo el criterio de clasificación y el formato esperado de la respuesta \cite{dong2024survey}. Las demostraciones, las instrucciones y la imagen que debe clasificarse forman el contexto disponible para generar esa respuesta. La Figura~\ref{fig:icl_context_related_work} ilustra el esquema general con una tarea de texto. Formalmente,
\begin{equation}
\mathcal{D}_K=\{(x_1,y_1),\ldots,(x_K,y_K)\},
\end{equation}
\end{samepage}
y predice la salida de \(x^\ast\) condicionada por esos ejemplos y la instrucción mediante
\begin{equation}
\hat{y}^\ast=f_\theta\left(x^\ast,\mathcal{D}_K,\text{instrucción}\right).
\end{equation}

Durante la consulta, los parámetros \(\theta\) permanecen fijos y la respuesta se adapta a los ejemplos presentes en el contexto. En cambio, el entrenamiento supervisado de la CNN incorpora los ejemplos mediante la actualización de sus pesos. \(K\) indica cuántas demostraciones recibe el VLM, aunque su selección, diversidad, orden y correspondencia entre imagen y etiqueta también afectan a la consulta. Los estudios sobre estos factores se recogen en la sección~\ref{sec:antecedentes_icl}.

\begin{figure}[H]
  \centering
  \safeincludegraphics[width=0.88\textwidth]{sources/icl_general_dong2024.pdf}
  \caption{Esquema de demostraciones y consulta en ICL, basado en Dong y colaboradores \cite{dong2024survey}.}
  \label{fig:icl_context_related_work}
\end{figure}
